% ED26-SRC-036-019-S19.5 referral [18]; captured author version April 2003.
\chapter{The Action Language C+: Laws, Events and Histories}
\label{sec:p16-cplus-action-language}

C+ separates what is caused from material constraints on a history.
Giunchiglia, Lee, Lifschitz, McCain and Turner use finite-domain constants
and causal laws to define a labelled transition system
\cite[Sections~4--5, Appendix~B]{GiunchigliaEtAl2003CplusAuthor}.
This treatment follows the openly offered author manuscript dated
8 April 2003, associated by the coauthor's publication page with the
2004 article; its pagination is not the publisher's pagination.
Its full three-class grammar is broader than the displayed fragment in
\cref{sec:p16-nct-cplus}. Later extensions are not silently included.
\index{C+}\index{action language}\index{fluent!statically determined}

\section{Signature and causal laws}
\label{sec:p16-cplus-grammar}

Use \cref{def:p16-nct-signature,def:p16-nct-theory}: each constant
has a nonempty finite domain, but the signature need not be finite.
An interpretation is explanatory when it uniquely satisfies the heads
whose bodies it makes true. The sign \(\Leftarrow\) constructs a causal
rule, not a material implication.
Constant-to-constant equality abbreviates
\[
 c=d\quad:\quad
 \bigvee_{v\in\operatorname{Dom}(c)\cap\operatorname{Dom}(d)}
                        (c=v\land d=v).
\]
An empty disjunction is \(\bot\); inequality negates this formula.
Boolean constants have values \(t,f\).

\begin{hbdefinition}{def:p16-cplus-laws}{The original three-class grammar}
Partition constants into fluents \(\sigma^{fl}\) and actions
\(\sigma^{act}\), and fluents into simple \(\sigma^{sim}\) and
statically determined \(\sigma^{sd}\).
Actions need not be Boolean.
A fluent formula contains only fluents, including constant-free formulas.
An action formula contains at least one action constant and no fluent.
A general formula may mix the two kinds.
An action description \(D\) is a set of laws:
\[
\begin{array}{lll}
 \text{static}&\text{caused }F\text{ if }G&
                         F,G\text{ fluent};\\
 \text{action dynamic}&\text{caused }F\text{ if }G&
                         F\text{ action},\ G\text{ general};\\
 \text{fluent dynamic}&\text{caused }F\text{ if }G\text{ after }H&
                         F,G\text{ fluent},\ H\text{ general}.
\end{array}
\]
A fluent dynamic head \(F\) contains no statically determined constant.
Heads may be compound formulas. Static and action dynamic laws share
keywords but have different head types.
\end{hbdefinition}

Static laws can explain derived values within a state. Action dynamic
laws can explain event attributes from fluents and other attributes.
Fluent dynamic laws connect origin state/event to successor fluents.

\begin{hbdefinition}{def:p16-cplus-translation}{The history translation}
For \(m\geq0\), use constants \(i:c\), with unchanged domains:
fluents at \(0\leq i\leq m\), actions at \(0\leq i<m\).
The notation \(i:F\) stamps every constant occurrence together.
Construct \(D_m\) by the following families:
\[
\begin{array}{ll}
 i:F\Leftarrow i:G&
  \text{static laws},\ 0\leq i\leq m;\\
 i:F\Leftarrow i:G&
  \text{action dynamic laws},\ 0\leq i<m;\\
 (i+1):F\Leftarrow((i+1):G)\land(i:H)&
  \text{fluent dynamic laws},\ 0\leq i<m;\\
 0:c=v\Leftarrow0:c=v&
  c\in\sigma^{sim},\ v\in\operatorname{Dom}(c).
\end{array}
\]
Only simple fluents receive initial exogeneity automatically.
Action exogeneity, inertia and derived-fluent exogeneity are not implicit.
\end{hbdefinition}

The \(G\) condition of a fluent dynamic law is tested at the successor;
\(H\) is tested at the origin, including its event. Replacing this by
an old-state precondition changes the law.

\section{All seventeen abbreviation conventions}
\label{sec:p16-cplus-abbreviations}

The author's Appendix B has seventeen numbered conventions, grouped
here without dropping any family. Every expansion must satisfy the
base grammar, including its dynamic-head restriction.

\begin{enumerate}
\item Omit \(\text{if }\top\) in a static or action dynamic law.
Write \(\text{caused }F\text{ after }H\) for a fluent dynamic law
with \(G=\top\).
\item \(\text{constraint }F\) means
\(\text{caused }\bot\text{ if }\neg F\), with \(F\) fluent.
\(\text{constraint }F\text{ after }G\) means
\(\text{caused }\bot\text{ if }\neg F\text{ after }G\),
with \(F\) fluent and \(G\) general.
These filter explanations rather than cause \(F\).
\item For any fluent \(c\), \(\text{rigid }c\) abbreviates
\(\text{constraint }c=v\text{ after }c=v\) for every value.
\(\text{always }F\) means
\(\text{caused }\bot\text{ after }\neg F\), for general \(F\).
It constrains origin-state/event combinations on transitions, not
all static states, and contributes nothing at horizon zero.
\item \(\text{nonexecutable }F\text{ if }G\), with action \(F\),
means \(\text{caused }\bot\text{ after }F\land G\).
For \(\text{F causes G if H}\), action formula \(F\) and fluent \(G\)
give \(\text{caused }G\text{ after }F\land H\);
action \(G\) instead gives \(\text{caused }G\text{ if }F\land H\).
The fluent head must contain only simple constants.
\item \(\text{default }F\text{ if }G\) means
\(\text{caused }F\text{ if }F\land G\), static or action dynamic
according to \(F\)'s type.
\(\text{default }F\text{ if }G\text{ after }H\) means
\(\text{caused }F\text{ if }F\land G\text{ after }H\);
its \(\text{if }G\) may be omitted when \(G=\top\).
\item \(\text{exogenous }c\text{ if }G\) means
\(\text{default }c=v\text{ if }G\) for all values.
For action \(F\), \(\text{F may cause G if H}\) means
\(\text{default }G\text{ after }F\land H\) when \(G\) is fluent,
and \(\text{default }G\text{ if }F\land H\) when \(G\) is action.
\item \(\text{inertial }c\text{ if }G\) means
\(\text{default }c=v\text{ after }c=v\land G\) for every value.
The fluent \(c\) must be simple; \(G\) is an origin condition.
A final \(\text{if }\top\) may be dropped in causes, default,
exogenous, may-cause and inertial declarations.
\item Appending \(\text{unless }a\) to static
\(\text{caused }F\text{ if }G\) requires Boolean derived \(a\)
and means the two static laws
\[
 \text{caused }F\text{ if }G\land\neg a,\qquad\text{default }\neg a.
\]
For an action dynamic law this requires Boolean action \(a\),
and both expanded laws are action dynamic.
For fluent dynamic \(\text{caused }F\text{ if }G\text{ after }H\),
it requires Boolean action \(a\) and means
\[
 \text{caused }F\text{ if }G\text{ after }H\land\neg a,\qquad
 \text{default }\neg a,
\]
the second law action dynamic. A state abnormality and an event
abnormality are thus different constant types.
\end{enumerate}

For instance inertia expands to
\[
 (i+1):c=v\Leftarrow
             ((i+1):c=v)\land(i:c=v)\land(i:G).
\]
The repeated successor atom licenses its persistence as an explanation
when the old value and condition hold. It is not a least-Horn-model rule.
Exogeneity supports false Boolean action values as well as true ones.
Non-Boolean actions have no distinguished no-action value unless the
description designates it.

\section{States, events and general paths}
\label{sec:p16-cplus-transition}

\begin{hbdefinition}{def:p16-cplus-system}{The labelled transition system}
A state \(s\) assigns every fluent and has \(0:s\) as an explanatory
model of \(D_0\). An event \(e\) assigns every action constant.
A transition \(\langle s,e,s'\rangle\) exists when
\(0:s\cup0:e\cup1:s'\) is an explanatory model of \(D_1\).
Its event label is a full assignment, not just a set of true actions.
A length-\(m\) path records \(s_0,e_0,s_1,\ldots,e_{m-1},s_m\);
a length-zero path is a state.
\end{hbdefinition}

\begin{hbtheorem}{thm:p16-cplus-paths}{proved=annex:p16.cplus-paths}{General state and path correspondence}
For every well-formed original C+ description, both endpoints of every
transition are states. For every finite \(m\geq0\), explanatory models
of \(D_m\) are exactly length-\(m\) paths, preserving all state/event
values. Arbitrary formula heads are allowed; neither definiteness nor
acyclicity is assumed.
\end{hbtheorem}

The proof partitions activated heads into coordinate blocks. It does
not apply an atom-head completion result to compound heads.

\begin{hbtheorem}{thm:p16-cplus-head-exception}{proved=annex:p16.cplus-head-exception}{The essential head restriction}
Let derived Boolean \(p\) have static law
\(\text{caused }p\text{ if }p\). Its only state has \(p=t\).
Translating the invalid law \(\text{caused }\neg p\text{ after }p\)
nevertheless gives explanatory history \(p_0=t,p_1=f\);
its target is not a state. This is an exception to a relaxed grammar,
not a legal C+ description.
\end{hbtheorem}

\begin{hbtheorem}{thm:p16-cplus-embedding}{proved=annex:p16.cplus-embedding}{Embedding causal theories}
For a causal theory \(T\), declare all constants derived fluents and
replace \(F\Leftarrow G\) by static \(\text{caused }F\text{ if }G\).
States are exactly \(T\)'s explanatory models.
With no further laws or actions, every ordered pair of states is a
transition labelled by the empty event. This describes the embedding,
not a physical dynamics inferred from a static theory.
\end{hbtheorem}

\section{Rigidity does not by itself license compression}
\label{sec:p16-cplus-rigid}

A fluent is rigid relative to \(D\) when every transition preserves its
value. For rigid set \(R\), form \(D_m^R\) by erasing timestamps on
its constants. Let \(I^R\) retain other coordinates and give each
\(c\in R\) its initial value.

The source illustrates rigidity in the monkey description by prohibiting
every \(PushBox(l)\). No remaining direct or static law explains a new
box location, so its inertia preserves its value on every transition.
The graph then separates into three box-location parts, with no edge
between different parts. This is a property of the modified description,
not a general consequence of calling a physical object fixed.

\begin{hbtheorem}{thm:p16-cplus-rigid}{proved=annex:p16.cplus-rigid}{Rigid-value compression with its hypotheses}
Suppose all members of \(R\) are derived fluents, and every law whose
head contains a member of \(R\) contains only members of \(R\)
throughout head and conditions. Then \(I\mapsto I^R\) bijects the
explanatory models of \(D_m\) and \(D_m^R\), for every finite \(m\).
Neither hypothesis can simply be removed.
\end{hbtheorem}

With simple Boolean \(p\), \(\text{rigid }p,\text{default }p\)
explain only \(p_0=p_1=t\) at horizon one, but compression also explains
\(p=f\). For derived \(p,q\), laws
\(\text{caused }p\text{ if }q,\text{exogenous }q\) have only state
\(p=q=t\), making \(p\) rigid. With \(R=\{p\}\), compression nevertheless
explains \(p=t,q_0=f,q_1=t\): a later cause explains the shared earlier
value. The annex gives both countermodels.

\section{The source's action descriptions}
\label{sec:p16-cplus-cases}

A description is definite when each head is a value atom or \(\bot\)
and no atom heads infinitely many laws. Then its \(D_m\) is definite,
so \cref{thm:p16-nct-completion} supplies an exact completion method.
Finite explicitly represented domains admit bounded SAT encoding.
No planning complexity or physical adequacy conclusion follows.

\subsection{A basic action}
Simple Boolean \(p\), Boolean action \(a\), and laws
\(a\text{ causes }p,\text{exogenous }a,\text{inertial }p\)
give \(p'=p\lor a\), including both initial and both event values.
False \(a\) preserves \(p\) because inertia was declared.

\subsection{The complete monkey-and-bananas law schemas}
Figures 2--3 declare simple \(\operatorname{Loc}(x)\), for
\(x\in\{Monkey,Box,Bananas\}\), with values \(L_1,L_2,L_3\), and simple
Boolean \(HasBananas,OnBox\).
Boolean actions are \(Walk(l),PushBox(l)\) for each location, and
\(ClimbOn,ClimbOff,GraspBananas\).
All actions are exogenous and all simple fluents inertial.
The static schemas are
\[
\begin{split}
 \text{caused }\operatorname{Loc}(Bananas)=l
   &\text{ if }HasBananas\land\operatorname{Loc}(Monkey)=l,\\
 \text{caused }\operatorname{Loc}(Monkey)=l
   &\text{ if }OnBox\land\operatorname{Loc}(Box)=l.
\end{split}
\]
\(Walk(l)\) causes the monkey to be at \(l\), and is nonexecutable if
already there or if on the box.
\(PushBox(l)\) causes both monkey and box to be at \(l\), and is
nonexecutable if the monkey is already there, is on the box, or is
not at the box.
\(ClimbOn\) causes \(OnBox\), and is prohibited if already on it or
not at it. \(ClimbOff\) causes \(\neg OnBox\), and is prohibited if
not on it. \(GraspBananas\) causes \(HasBananas\), and is prohibited
if already possessing them, not on the box, or not at the bananas.
For every \(l\), prohibit \(Walk(l)\land PushBox(l)\),
\(Walk(l)\land ClimbOn\), \(PushBox(l)\land ClimbOn\);
also prohibit \(ClimbOff\land GraspBananas\).
This is not a ban on every concurrent action combination.

From monkey and box at \(L_1\), bananas at \(L_3\), both Booleans false,
the successive events \(PushBox(L_3),ClimbOn,GraspBananas\), with all
other actions false each time, give a three-step path.
Carrying the bananas uses the static ramification.

Figure 4 extends the signature to a fixed set \(M\) of monkey names:
\(\alpha,\beta\) range over \(M\), location constants are indexed by
\(M\cup\{Bananas,Box\}\), and each monkey has its own Boolean
\(HasBananas(\alpha),OnBox(\alpha)\) and indexed action constants.
Its David and Goliath exceptions concern designated members of that set.
The multi-monkey description uses derived Boolean
\(Ab_1(\alpha,\beta)\) on the static occupancy constraint
\(\text{constraint }\neg(OnBox(\alpha)\land OnBox(\beta))
\text{ unless }Ab_1(\alpha,\beta)\), for \(\alpha\ne\beta\).
Static causes make \(Ab_1(David,\alpha)\) and \(Ab_1(\alpha,David)\)
true for its exception.
Boolean action \(Ab_2(\alpha)\) qualifies both effects of
\(PushBox(\alpha,l)\). Action dynamic laws cause it when
\(OnBox(Goliath)\) or \(ClimbOn(Goliath)\).
Negative defaults come from unless; abnormalities are not declared
exogenous. Defeating push effects is different from prohibiting push.
These are the displayed abnormality mechanisms, not a claim to have
enumerated every multi-monkey transition.

\subsection{Conditional shooting}
Figure 5 has simple Boolean \(Loaded,Alive(Turkey_1),Alive(Turkey_2)\),
simple \(Target\in\{Turkey_1,Turkey_2,None\}\), and Boolean actions
\(Load,Aim(Turkey_1),Aim(Turkey_2),Shoot\).
Actions are exogenous and fluents inertial.
Load causes \(Loaded\) and \(Target=None\); \(Aim(x)\) causes \(Target=x\).
Shoot causes \(\neg Alive(x)\) if \(Target=x\), and causes \(\neg Loaded\).
Prohibit shooting unloaded and each \(Aim(x)\land Shoot\).
The target condition is tested before shooting.
With target \(None\), shooting unloads but kills neither turkey.

\subsection{Nondeterministic travel}
Figure 6 has simple locations of Jack and car in \(\{Home,Work\}\)
and exogenous Boolean \(Go(Home),Go(Work)\); both locations are inertial.
\(Go(l)\) causes \(\operatorname{Loc}(Jack)=l\) and may cause
\(\operatorname{Loc}(Car)=l\) if
\(\operatorname{Loc}(Car)=\operatorname{Loc}(Jack)\).
It is prohibited if Jack is already at \(l\).
Optional explanation of the car's new value is not a probability law.

\subsection{Concurrent lifting and spontaneous tendency}
Figure 7 has simple \(Level(LeftEnd),Level(RightEnd)\in\{Low,High\}\)
and Boolean \(OnTable\), all inertial, and exogenous Boolean \(Lift(x)\).
A lift causes its end high and is prohibited if already high.
The static law causes \(\neg OnTable\) when levels differ.
Starting low at both ends and on the table, one lift takes the object
off; two simultaneous lifts leave it on.

Figure 8 has simple Boolean \(Right\), not declared inertial,
and exogenous Boolean \(Hold\).
Hold causes \(Right\) if \(Right\), and \(\neg Right\) if \(\neg Right\).
Defaults \(Right\text{ after }\neg Right\) and
\(\neg Right\text{ after }Right\) supply a tendency to flip.
Persistence when held is an action effect, not ordinary inertia.

\begin{hbtheorem}{thm:p16-cplus-outcomes}{proved=annex:p16.cplus-outcomes}{Selected exact outcomes}
The basic action has its four \(p'=p\lor a\) transitions.
The stated three-step monkey history is a path.
From a loaded state targeting \(Turkey_1\), both turkeys alive, a
shooting-only event unloads and leaves only \(Turkey_2\) alive.
From Jack and car at Home, the Go(Work)-only event has exactly two
successors: Jack at Work and car at Home or Work.
The one-lift and two-lift events have the stated unique successors.
The pendulum preserves Right exactly when held, otherwise flips.
\end{hbtheorem}

\subsection{Definitions and non-Boolean event attributes}
Section 5.5 adds derived Boolean \(NextToBox\) with
\[
 \text{caused }NextToBox\text{ if }
    \operatorname{Loc}(Monkey)=\operatorname{Loc}(Box),\qquad
 \text{default }\neg NextToBox.
\]
It can replace the equality in conditions. A simple declaration would
add initial self-support and could make it incorrectly true initially.

Figure 9 declares simple inertial Boolean \(HasPublications,
HasJournalPublications,HasLongPublications\).
Its exogenous action constants are Boolean \(Publish\),
\(Kind\in\{Conference,Journal,None\}\) and
\(Length\in\{1,\ldots,100,None\}\).
The laws are
\[
\begin{gathered}
 Publish\text{ causes }HasPublications,\\
 \text{always }((Kind=None)\equiv\neg Publish),\\
 Publish\text{ causes }HasJournalPublications\text{ if }Kind=Journal,\\
 \text{always }((Length=None)\equiv\neg Publish),\\
 Publish\text{ causes }HasLongPublications\text{ if }Length=l
                                      \quad(31\leq l\leq100).
\end{gathered}
\]
The always-laws constrain event attributes, not static fluent values.
The source first contrasts this with changing the action vocabulary from
\(Publish\) to \(Publish(k)\), and then to \(Publish(k,l)\): adding a
kind and a length argument requires rewriting the old effect schemas.
The attribute representation instead retains the action constant and
adds event coordinates and their constraints. This is the stated
representation tradeoff, not an assertion that every parameterized
action notation is inferior.

\begin{hbtheorem}{thm:p16-cplus-definitions}{proved=annex:p16.cplus-definitions}{Definitions and attributes}
For fresh derived Boolean \(d\) and old-signature fluent formula \(F\),
adding \(\text{caused }d\text{ if }F,\text{default }\neg d\)
uniquely extends each state and history by \(d\equiv F\).
Replacing occurrences of \(F\) in old law conditions by \(d\) preserves
the correspondence, keeping old heads unchanged.
The simple-fluent variant fails already at horizon zero.
The publication description has one nonpublishing event, both attributes
None, and \(200\) publishing events.
From an all-false state, Journal with length 31 makes all three fluents
true; Conference with length 20 makes only HasPublications true.
Each specified event has a unique successor.
\end{hbtheorem}

\section{Scope of the causal representation}
\label{sec:p16-cplus-limits}

The correspondence is internal to the given laws and finite horizon.
Constraints can prohibit events; incompatible caused values can imply
additional prohibitions. Defeated effects need not prohibit actions.
May-cause laws permit multiple explained successors without probabilities.
Derived definitions, inertia and default tendencies provide different
kinds of support.

An event label is not automatically a structural-model intervention,
nor does an explanatory model identify a mechanism from observations.
Rules stipulate what counts as an explanation, without establishing the
physical reliability of these idealizations. Stable-model,
situation-calculus and temporal-action encodings require separate
translation hypotheses. The captured source's later CCalc discussion
and subsequent language extensions are not used to assert a broader
algorithm or correspondence.
